\documentclass{article}
\usepackage[utf8]{inputenc}
\usepackage{amsmath}
\usepackage{geometry}
\geometry{margin=2cm}
\begin{document}

\section*{Questão 142 — ENEM 2024 — Matemática}

Uma criança, utilizando um aplicativo, escreveu uma mensagem para um amigo. Esta mensagem foi produzida seguindo etapas descritas abaixo:

\begin{itemize}
  \item Na 1\textordfeminine{} etapa, inseriu três figuras do tipo 🙂 no visor da mensagem.
  \item Na 2\textordfeminine{} etapa, copiou o que havia escrito na etapa anterior e colou ao lado, duplicando as figuras.
  \item Na 3\textordfeminine{} etapa, repetiu o processo, copiando e colando o conteúdo da etapa imediatamente anterior, e assim sucessivamente até a 20\textordfeminine{} etapa.
\end{itemize}

\noindent \textbf{Pergunta:} Qual o total de figuras contidas na mensagem enviada na 20\textordfeminine{} etapa?

\section*{Resolução Técnica}

A quantidade de figuras dobra a cada etapa, iniciando com 3 na primeira. Para determinar o total na etapa $n$, usa-se a fórmula da progressão geométrica:

\[
a_n = 3 \times 2^{n-1}.
\]

Assim, para a 20\textordfeminine{} etapa:

\[
a_{20} = 3 \times 2^{19}.
\]

\section*{Explicação Didática}

O processo inicia com 3 figuras, que são duplicadas a cada nova etapa, configurando uma progressão geométrica de razão 2. Esta sequência é representada por $a_n = 3 \times 2^{n-1}$. O cálculo para a etapa 20 gera o resultado final pedido.

\section*{Erros Comuns e Distratores}

Algumas alternativas incorretas podem surgir por:

\begin{itemize}
  \item Usar expoente 20 em vez de 19.
  \item Aplicar alterações sem base, como subtrações indevidas.
  \item Confundir índices de expoentes.
\end{itemize}

\section*{Perfil da Questão}

Esta questão avalia a habilidade de reconhecer e aplicar progressão geométrica, com demanda cognitiva média e necessidade de interpretar sequências numéricas e manipular expoentes.

\section*{Tabela do Enunciado (Imagem 1 Descritiva)}

\begin{tabular}{|p{7cm}|p{7cm}|}
\hline
\textbf{Etapas} & \textbf{Visor de escrita} \\
\hline
1\textordfeminine{} etapa: inseriu três figuras 🙂 no visor da mensagem; & 🙂 🙂 🙂 \\
\hline
2\textordfeminine{} etapa: copiou o que havia inserido anteriormente e colou ao lado; & 🙂 🙂 🙂 🙂 🙂 🙂 \\
\hline
3\textordfeminine{} etapa: copiou o que estava no visor na 2\textordfeminine{} etapa e colou ao lado; & 🙂 🙂 🙂 🙂 🙂 🙂 🙂 🙂 🙂 🙂 🙂 🙂 \\
\hline
\end{tabular}

\bigskip

Este padrão mostra claramente que a cada etapa, o número de figuras dobra, confirmando o modelo matemático apresentado.


\end{document}